\section{Introduction}\label{sec:intro}

As the world’s largest athletics association, the Deutscher Leichtathletik-Verband (DLV) is tasked with facilitating peak individual performance and maintaining an internationally 
competitive national team \citep{DLV2026}. However, zero medals at the 2023 World Athletics Championships in Budapest and a historic low medal count at the 2024 Paris Olympics have 
left the organization in a sobering position. Although individual athletes still occasionally reach the podium, the declining return on investment in terms of Olympic success \citep{fremerey2024} 
raises questions about the underlying health of the national talent pool.

While medals are the most visible metric of success, they reflect only the peak of elite performance and offer little insight into how performance is distributed within the elite level. 
To examine potential structural weaknesses, we analyze the top 35 rankings from 2001 to 2025, guided by two research questions. First, we assess the impact of the COVID-19 pandemic as a 
short-term disruption to the development of elite performances. 
Secondly, we are investigating whether the performance gap between the national elite and the broader field of competitors has widened and how performance has generally developed in comparison. 
This investigation is prompted by observations at national championships where a growing decline in performance was noted beyond the top positions. 
We aim to quantify this reduction in competitive intensity, as a lack of internal pressure, in line with the principle that "competition fosters performance," can lead to top athletes falling behind internationally.

Our analytical approach is detailed in \Cref{sec:methods}, where we describe the data processing and justify the analytical framework. Building on this, we will analyze performance metrics 
to identify structural trends and evaluate our research questions (\Cref{sec:results}). Finally, \Cref{sec:conclusion} discusses the implications of these findings.

\section{Data and Methods}\label{sec:methods}
\subsection{Data Acquisition and Standardization}           
 
The data were retrieved from the official archive by the Deutscher Leichtathletik-Verband (DLV), which documents the top performances across the major athletic disciplines (including running 
events, hurdles and steeplechase, jumping events, throwing events, and combined events) from 2001 to 2025. The archive covers performances across gender (male, female) and age groups (under 14 up to adults). 

As the data for each year, age group, and gender are stored in separate PDF files, we used \textit{pdfplumber} in combination with specific regular expressions to handle formatting differences 
over years. Outliers due to entry inconsistencies in the original files (e.g., irregular time formats) were corrected when the intended values could be inferred from external sources. Minor data 
gaps resulting from parsing errors were manually completed by referencing the original PDF files. Non-recurrent distances (e.g., [3000 m steeplechase and] 7.5 km events in selected youth categories) 
were excluded, yielding a final dataset of 226,683 entries.

We standardized discipline names and birth years to a four-digit format to ensure consistency across the dataset. Performance marks were converted into numeric values (seconds for track events and 
meters for field events) to allow statistical analyses. To facilitate comparisons across events, performances were transformed using World Athletics scoring tables, which express results on a common 
point scale. To calculate the points for each performace, we used XXX